%=============================================================================
\section{研究方法：人机协作的理论物理研究}
\label{sec:methodology}
%=============================================================================

CTUFT 的发展代表了理论物理研究的一种新范式，即将人类物理直觉与
人工智能辅助的计算和验证相结合。本节概述本工作中采用的协作研究方法。

\subsection{分工}

研究工作流围绕"人类主导、AI 辅助"的框架展开：

\begin{itemize}
    \item \textbf{人类（主导）}：负责概念框架的构建、数学结构的物理诠释、
          可证伪预测的设计，以及整体理论一致性把控；
    
    \item \textbf{AI（辅助）}：负责复杂的代数运算、大规模数值验证、
          自动化代码生成，以及全面的文献综合。
\end{itemize}

这种分工确保核心物理洞察和诠释基础始终扎根于人类的科学推理，
同时利用 AI 的计算可扩展性和快速迭代能力。

\subsection{验证协议}

为确保可靠性和科学严谨性，所有 AI 生成的计算和推导都经过
系统化的验证协议：

\begin{enumerate}
    \item \textbf{交叉验证}：AI 生成的结果通过与独立计算方法
          和解析近似（如可用）进行核对；
    
    \item \textbf{物理一致性}：所有导出表达式必须满足已知的物理约束，
          包括量纲分析、对称性要求和守恒律；
    
    \item \textbf{极限情况分析}：通过解析可处理的极限情况测试结果，
          以验证渐近行为和边界条件。
\end{enumerate}

这一验证框架确保，尽管 AI 加速了研究进程——使得复杂数学结构的
快速探索成为可能——但物理内容和理论结论始终牢固地锚定于
成熟的科学方法和人类监督之中。

\subsection{对理论物理的启示}

本工作展示的人机协作方法，暗示了未来理论物理研究的一种变革性模式。
通过将人类的创造力和物理洞察与 AI 处理复杂计算和模式识别的能力相结合，
研究人员可以探索通过传统方法难以处理的参数空间和数学结构。
这种范式并非取代人类判断，而是放大它，有可能加速那些既需要
深刻理解概念又需要大量计算资源的领域的发现步伐。
